\section{Selecting the final application by complete time}
\label{sec:notes-fast-deployment}
We selected the new tree using recurrence times. A lower kernel cost alone does not establish a lower application cost, so deployment requires a further check on the original validation split. We run the same 56 validation requests in fresh processes, with prediction, exact reduction, and coefficient comparison inside each timer. The two choices use the same fast implementation and retain the same native seed model; they differ only in the recurrence strategy. Three paired repetitions provide each request's median time. We enable the new recurrence model only when its summed median time is at least two percent lower and its paired repetition interval stays below one. The 65 reserved requests do not enter this decision.

The complete validation total is 13.282 seconds with the recurrence model and 13.408 seconds with the fixed strategy. The ratio is 0.991, with repetition interval $[0.953,1.015]$. The final application therefore uses the fixed recurrence strategy. The kernel saving does not meet the complete-time deployment rule, so the trained tree remains a recorded candidate. The retained sparse-seed model remains active in its validated coupled domain. This decision leaves all exact formulas and coefficient checks unchanged.

A separate twelve-request test follows the frozen decision. It includes six rational polygons, two distinct-triangle products at two and thirteen loops, a nonplanar box, a symbolic bubble, and two thirteen-loop analytic families. Appendix~\ref{app:notes-fast-deployment} specifies their fresh inputs. Table~\ref{tab:notes-fast-deployment} keeps this test separate from the preceding 65-request experiment.
\begin{table}[!htbp]\centering\small
\caption{Separate reserved test of the final complete-cost rule. Times are sums of medians in seconds. ``Candidate'' always enables the recurrence model. ``Fixed recurrence'' retains the native seed model but uses the fixed recurrence strategy.}
\label{tab:notes-fast-deployment}
\begin{tabularx}{\linewidth}{@{}Yrrrrr@{}}\toprule
Collection & $N$ & Final & Candidate & Fixed recurrence & Previous\\\midrule
Scalar recurrence & 6 & 1.708 & 1.670 & 1.756 & 4.125 \\
Distinct products & 2 & 0.614 & 0.600 & 0.598 & 0.791 \\
Nonplanar box & 1 & 3.856 & 3.719 & 3.818 & 3.955 \\
Symbolic bubble & 1 & 0.699 & 0.737 & 0.724 & 0.704 \\
Thirteen-loop formulas & 2 & 0.377 & 0.323 & 0.351 & 1.763 \\
All requests & 12 & 7.254 & 7.048 & 7.248 & 11.338 \\
\bottomrule\end{tabularx}\end{table}

Against the preceding deployed application, the final method uses 36.0\% less time, with ratio 0.640 and repetition interval $[0.612,0.712]$. Against the always-enabled recurrence candidate, the final method uses 2.9\% more time, with ratio 1.029 and repetition interval $[0.981,1.131]$. This interval includes equal times. Against the fixed recurrence strategy, the final method uses 0.1\% more time, with ratio 1.001 and repetition interval $[0.951,1.118]$. This interval includes equal times.

The final method completes 12 of twelve requests in every repetition. All verified outputs agree with the independent references, with 2,856 coefficient comparisons in this test and 6,072 in the full-cost validation. Reference preparation for the twelve requests takes 73.895 seconds and remains separate. The broad audit in the preceding section belongs to the initial candidate. The final test checks the deployment rule on fresh inputs, while the earlier audit retains its original outcomes and native limits.
